Definiu-se que existirão três tipos de utilizadores: o administrador, o docente e o aluno/concorrente.\\

\begin{itemize}
  \item Administrador - é a entidade com mais poder no sistema. É o único que pode criar contas do tipo docente. É caracterizado por:
    \begin{itemize}
      \item Nome de utilizador;
      \item Nome completo;
      \item Password
      \item e-Mail
    \end{itemize}
  \item Docente - entidade que tem permissões para criar concursos, exercícios, aceder aos resultados das submissões, (...). Os seus atributos coincidem com os do Administrador.
  \item Aluno/Concorrente - (...)
\end{itemize}

Decidiu-se que o sistema terá a noção de grupo. 
(...)
O grupo é que possui as credenciais para entrar no sistema (nome de utilizador e password). 
Além disso também tem um nome pelo qual é identificado, um e-mail que será usado no caso de haver necessidade de se entrar em contacto com o grupo, e um conjunto de concorrentes.\\

\begin{figure}[htbp]
\begin{center}
\includegraphics[width=0.5\textwidth]{Images/missimage}
\caption{Modelo de dados}\label{fig modedados}
\end{center}
\end{figure}

Achamos importante incluir a informação de cada concorrente no grupo para, se possível, automatizar várias tarefas tais como lançamento de notas.
Cada concorrente é caracterizado pelo seu nome completo, número de aluno (se for o caso), e e-mail.\\

Para finalizar vamos explicar em que consistem os concursos, exercícios e tentativas.

Um concurso, resumidamente, é um agregado de enunciados. 
Tem outras propriedades tais como um título, data de inicio e data de fim (período em que o concurso está disponível para que os grupos se registem), 
chave de acesso (necessária para o registo dos grupos), duração do concurso (tempo que o grupo tem para resolver os exercícios do concurso, 
a partir do momento que dá inicio à prova), e por fim, regras de classificação.\\

Um enunciado é um exercício que o concorrente tenta solucionar. Como seria de esperar, cada exercício pode ter uma cotação diferente, 
logo o peso do enunciado é guardado no mesmo. 
Para cada enunciado existe também um conjunto de inputs e outputs, que servirão para verificar se o programa submetido está correcto. 
Contém ainda uma descrição do problema que o concorrente deve resolver, assim como uma função de avaliação, função esta que define como se verifica se o output obtido está de acordo com o esperado.\\
